\documentclass[11pt]{article}

\usepackage[margin=1in]{geometry}
\usepackage{amsmath,amssymb}
\usepackage{booktabs}
\usepackage{hyperref}

\title{Seed-Universal Disorder Flux Release for HAOS-IIP}
\author{Tomislav Rupi\'c \\ Independent Researcher}
\date{April 2026}

\begin{document}
\maketitle

\begin{abstract}
This paper freezes release \texttt{53.4} of HAOS-IIP as the closure of the remaining seed-universal disorder-native flux margin. The earlier disorder-native flux result already passed on aggregated observables, but one tested seed in the hardest smooth-radial mild-disorder case left the bridge row at \texttt{WATCH}, with only \texttt{2/3} seeds passing. Release \texttt{53.4} keeps the same scalar carrier, the same native-flux reconstruction, and the same thresholds, but extends the delayed asymptotic fit window from \([0.006, 0.020]\) to \([0.006, 0.024]\). Under that bounded readout, all twelve aggregated mild-disorder cases still pass, the maximum refinement drift improves to \texttt{0.1249}, and the minimum seed-level pass count becomes \texttt{3/3}. The correct conclusion is therefore bounded and external: tested seed-universal disorder-native flux closure now holds on the smooth-radial mild-disorder window, without modifying HAOS core or promoting a universality theorem.
\end{abstract}

\tableofcontents

\section{Scope}

Release \texttt{53.4} addresses a narrow remaining bridge boundary:
\begin{quote}
can the disorder-native flux proxy close at the seed-universal level on the tested smooth-radial mild-disorder window?
\end{quote}

This release does not change HAOS core, alter the scalar carrier, or loosen thresholds. It changes only the external fit window used by the native-flux runner.

\section{Construction}

The carrier and transport construction remain fixed:
\begin{itemize}
\item same scalar kernel-graph family;
\item same smooth radial deformation amplitudes \(\eta = 0.05, 0.10, 0.15\);
\item same mild disorder fractions \(\sigma = 0.02, 0.04\);
\item same tested seeds \(\{0,1,2\}\);
\item same native bulk shells, interior-ball cumulative flux, and local-linear radial density gradient;
\item same thresholds for median error, p90 error, diffusivity gap, shell-\(\kappa\) CV, refinement drift, and metric tracking.
\end{itemize}

The only change is the delayed asymptotic fit window:
\begin{itemize}
\item previous fit window: \([0.006, 0.020]\);
\item updated fit window: \([0.006, 0.024]\).
\end{itemize}

\section{Frozen Result}

\begin{table}[h!]
\centering
\begin{tabular}{@{}lccccc@{}}
\toprule
Quantity & Value \\
\midrule
Aggregated passing cases & 12/12 \\
Minimum seed-level pass count & 3/3 \\
Maximum refinement drift & 0.1249 \\
Weakest aggregated case & \texttt{radial\_native\_flux|n=13|eta=0.150|sigma=0.020} \\
Hardest seed-level case & \texttt{radial\_native\_flux|n=13|eta=0.150|sigma=0.020|seed=1} \\
\bottomrule
\end{tabular}
\caption{Seed-universal disorder-native flux closure on the tested window.}
\end{table}

The weakest aggregated case has:
\begin{itemize}
\item \(\kappa / D_{\mathrm{eff}} = 0.9737\);
\item median relative error \texttt{0.0821};
\item p90 relative error \texttt{0.2876};
\item shell-\(\kappa\) CV \texttt{0.1177};
\item \(|\mathrm{metric\ tracking\ corr}| = 0.9944\).
\end{itemize}

The hardest seed-level trial now has:
\begin{itemize}
\item median relative error \texttt{0.0895};
\item p90 relative error \texttt{0.2660};
\item shell-\(\kappa\) CV \texttt{0.1092};
\item pass state \texttt{TRUE}.
\end{itemize}

\section{Interpretation}

The old bridge statement was:
\begin{quote}
aggregate disorder-native transport passes, but seed-universal closure remains watched at \texttt{2/3}.
\end{quote}

The updated statement is:
\begin{quote}
the same disorder-native flux reconstruction now closes at the tested seed-universal level on the smooth-radial mild-disorder window.
\end{quote}

This is still a bounded scalar-carrier transport result. It does not establish arbitrary disorder universality, a conservation theorem, or continuum ontology.

\section{Non-Claims}

This release does not claim:
\begin{itemize}
\item arbitrary disorder universality;
\item strong-disorder closure;
\item current conservation theorems;
\item curvature extraction;
\item continuum field equations;
\item equivalence to established physical theories.
\end{itemize}

\section{Artifacts}

\begin{itemize}
\item script: \path{numerics/simulations/scalar_kernel_graph_current_closure_radial_disorder_native_flux.py};
\item operator note: \path{experiments/operators/Scalar_Kernel_Graph_Current_Closure_Radial_Disorder_Native_Flux_v1.md};
\item frozen result: \path{data/20260423_184421_scalar_kernel_graph_current_closure_radial_disorder_native_flux.json};
\item latest result: \path{data/scalar_kernel_graph_current_closure_radial_disorder_native_flux_latest.json};
\item summary plot: \path{plots/20260423_184421_scalar_kernel_graph_current_closure_radial_disorder_native_flux_summary.png};
\item profile plot: \path{plots/20260423_184421_scalar_kernel_graph_current_closure_radial_disorder_native_flux_profiles.png}.
\end{itemize}

\section{Next Boundary}

The next honest step is to widen the seed set or add one stronger disorder fraction before using broader disorder-robust language.

\section{Conclusion}

Release \texttt{53.4} promotes the disorder-native flux seed margin from \texttt{WATCH} to bounded \texttt{PASS}. The tested smooth-radial mild-disorder window now closes at \texttt{3/3} seeds under the same native-flux reconstruction and the same thresholds, with only a later asymptotic readout window. That is a stronger bounded transport statement, not a universality theorem.

\end{document}
